% ECONOMICS_PROBLEM: {"id": "C65-33", "title": "Integrable terminal data at the critical logarithmic BSDE growth rate", "jel_code": "C65", "source_id": "OP-0562"}
% FIRST_POSTED: 2026-10-10
% ATTEMPT_STATUS: solved
% ENTRY_KIND: research_attempt
% AUTHOR: Ji Ho Bae
% AI_ASSISTANCE: GPT Codex; independent adversarial mathematical checks; no Lean claim
% SOURCE_REPOSITORY: https://github.com/jbaelaw/critical-log-bsde-counterexample
\documentclass[11pt]{article}
\usepackage[a4paper,margin=25mm]{geometry}
\usepackage[T1]{fontenc}
\usepackage{lmodern,microtype,amsmath,amssymb,amsthm}
\usepackage[colorlinks=true,linkcolor=blue,citecolor=blue,urlcolor=blue]{hyperref}
\newtheorem{theorem}{Theorem}
\newtheorem{lemma}{Lemma}
\newcommand{\E}{\mathbb E}
\newcommand{\F}{\mathcal F}
\newcommand{\R}{\mathbb R}
\title{An $L^1$ counterexample at critical logarithmic BSDE growth}
\author{Ji Ho Bae}
\date{10 October 2026}
\begin{document}
\maketitle
\begin{abstract}
For the generator $g(z)=|z|/\sqrt{\log(e+|z|)}$, we construct a nonnegative
terminal value of expectation one for which the scalar BSDE has no
class-$(D)$ solution. The construction works on every finite positive
horizon and in every Brownian dimension. It gives a negative answer to
Fan, Hu and Tang's Problem~5.1, with the original generator and solution class.
\end{abstract}

Let $B$ be a $d$-dimensional Brownian motion, $d\geq1$, on the Brownian
basis of \cite{FHT}, and fix $T>0$. Class $(D)$ means that
$\{Y_\tau:\tau\leq T\text{ a stopping time}\}$ is uniformly integrable.
A solution has continuous adapted $Y$, predictable $Z$, and
$\int_0^T|Z_t|^2dt<\infty$ almost surely, with the BSDE identity below.
All logarithms are natural; vector norms are Euclidean.

\begin{theorem}\label{main}
There is an explicitly specified $\F_T$-measurable $\xi\geq0$ with
$\E\xi=1$ for which no class-$(D)$ solution satisfies
\begin{equation}\label{bsde}
 Y_t=\xi+\int_t^T g(Z_s)\,ds-\int_t^T Z_s\cdot dB_s,
 \qquad g(z)=\frac{|z|}{\sqrt{\log(e+|z|)}}.
\end{equation}
\end{theorem}

We first record why bounded lower solutions can be compared with a
putative class-$(D)$ solution despite a change of measure.

\begin{lemma}\label{comparison}
If $(Y,Z)$ solves \eqref{bsde} with $\xi\geq0$ and $Y$ of class $(D)$,
then $Y\geq0$. Suppose a continuous adapted process $u$ satisfies
$0\leq u\leq C$ for a deterministic finite constant $C$, and
\[
 du_t=-f_t\,dt+z_t\cdot dB_t,
 \quad \int_0^T|z_t|^2dt<\infty,
 \quad f_t\leq g(z_t),\quad u_T\leq\xi.
\]
Then $u\leq Y$.
\end{lemma}
\begin{proof}
Since $g\geq0$, localization of \eqref{bsde} gives
$Y_t\geq\E[Y_{\tau_n}\mid\F_t]$ for stopping times $\tau_n\uparrow T$.
The stopped stochastic integrals may be chosen square-integrable.
Class $(D)$ and continuity give $Y_{\tau_n}\to\xi$ in $L^1$, whence
$Y_t\geq\E[\xi\mid\F_t]\geq0$. Taking rational $t$ and using continuity
makes this assertion simultaneous in time.

For $h(r)=r/\sqrt{\log(e+r)}$, $0\leq h'(r)\leq1$ on $r>0$.
Thus $g$ is globally $1$-Lipschitz, also in $\R^d$.
Put $\Delta=u-Y$ and choose the predictable vector
\[
 \theta_t=\begin{cases}
 \displaystyle\frac{g(z_t)-g(Z_t)}{|z_t-Z_t|^2}(z_t-Z_t),&z_t\ne Z_t,\\
 0,&z_t=Z_t.
 \end{cases}
\]
Then $|\theta_t|\leq1$. Novikov's condition holds, so under the equivalent
measure with density
$\exp(\int_0^T\theta_t\cdot dB_t-\frac12\int_0^T|\theta_t|^2dt)$,
the process $B^{\theta}=B-\int_0^{\cdot}\theta_tdt$ is Brownian. Writing
$r_t=g(z_t)-f_t\geq0$, we obtain
\[
 d\Delta_t=r_tdt+(z_t-Z_t)\cdot dB_t^{\theta}.
\]
Tanaka's formula makes $\Delta^+$ a local submartingale. It is bounded
by $C$, because $Y\geq0$, hence is a true submartingale. Its terminal
value is zero, so it vanishes identically. This argument needs no
integrability of $Y$ or $\xi$ under the changed measure.
\end{proof}

\section*{Bounded Gaussian lower solutions}
Write $\Phi$ and $\phi$ for the standard normal distribution function
and density. Set
\[
 S=\min\{T,(2\pi)^{-1}\},\qquad
 H_N=\sum_{k=1}^N\frac1k,\qquad
 \delta_k=\frac{S}{kH_N},\qquad t_k=\sum_{j=1}^k\delta_j.
\]
For each integer $N\geq1$, define
\begin{equation}\label{terminalN}
 X_N=2^N\prod_{k=1}^N
 \Phi\left(\frac{B^1_{t_k}-B^1_{t_{k-1}}}{\sqrt{\delta_k}}\right).
\end{equation}
The standardized increments are independent standard normals, and
$\E\Phi(G)=1/2$ by symmetry. Consequently $\E X_N=1$.

Define deterministic constants backwards, starting with $k=N$:
\begin{equation}\label{backwards}
 D_k=(2\pi\delta_k)^{-1/2},\quad
 L_k=2^ND_k\prod_{j>k}p_j,\quad
 q_k=[\log(e+L_k)]^{-1/2},\quad
 p_k=\Phi(q_k\sqrt{\delta_k}/\sqrt2).
\end{equation}
This recursion uses only constants already defined. For
$t_{k-1}\leq t\leq t_k$, put
\[
 A_k=2^N\prod_{j<k}\Phi(\Delta B^1_j/\sqrt{\delta_j})
                  \prod_{j>k}p_j,
 \qquad
 w_t=\frac{B^1_t-B^1_{t_{k-1}}+q_k(t_k-t)}{\sqrt{\delta_k+t_k-t}},
\]
where $\Delta B^1_j=B^1_{t_j}-B^1_{t_{j-1}}$, and define
\[
 u^N_t=A_k\Phi(w_t),\qquad
 z^N_t=\frac{A_k\phi(w_t)}{\sqrt{\delta_k+t_k-t}}\,e_1.
\]
At $t_k$, the next stage starts with its Gaussian factor equal to
$p_{k+1}$, so the values of $u^N$ agree. After $S$, set
$u^N=X_N$ and $z^N=0$. Direct It\^o differentiation on each interval gives
\begin{equation}\label{ito}
 du^N_t=-q_k|z^N_t|dt+z^N_t\cdot dB_t,\qquad
 0\leq |z^N_t|\leq L_k.
\end{equation}
The finite collection of deterministic bounds gives square-integrability
of $z^N$. Moreover $0\leq u^N\leq2^N$, $u^N_T=X_N$, and
\[
 q_k|z^N_t|\leq g(z^N_t),\qquad u^N_0=\prod_{k=1}^N2p_k.
\]
Thus $u^N$ satisfies the bounded lower-solution conditions of
Lemma~\ref{comparison}.

\section*{Divergence of the initial values}
Put $x_k=q_k\sqrt{\delta_k}/\sqrt2$ and $\ell_k=\log(2p_k)$.
Since $D_k\geq1$ and $p_j\geq1/2$, we have $L_k\geq2^k$ and
$q_k\leq(k\log2)^{-1/2}$. Also $q_k\leq1$, so $0\leq x_k\leq1$.
On this interval,
\begin{equation}\label{normalbound}
 \phi(1)x\leq\log(2\Phi(x))\leq\sqrt{2/\pi}\,x.
\end{equation}
Indeed the derivative is $\phi(x)/\Phi(x)$, between $\phi(1)$
and $2\phi(0)$. It follows that
\begin{equation}\label{tail}
 \sum_{j>k}\ell_j
 \leq\sqrt{\frac{S}{\pi\log2}}\,\sqrt{H_N}.
\end{equation}
Set $a=\log(1/(2\pi S))\geq0$. For $k\geq\lceil H_N\rceil$ and
$H_N\geq\max\{8,a\}$, use
\[
 \log L_k=k\log2+\tfrac12\log(kH_N)+\tfrac12a+\sum_{j>k}\ell_j.
\]
The coefficient in \eqref{tail} is at most one. Since $L_k\geq2^k$,
$\log(e+L_k)\leq1+\log L_k$, and hence
\[
 \log(e+L_k)\leq1+k\log2+\log k+k/2+\sqrt{k}\leq5k.
\]
The lower bound in \eqref{normalbound} now gives
\begin{align}
 \log u^N_0
 &\geq\phi(1)\sqrt{\frac{S}{10H_N}}
            \sum_{k=\lceil H_N\rceil}^{N}\frac1k\notag\\
 &\geq\frac{\phi(1)}{2\sqrt{10}}\sqrt{SH_N}.\label{divergence}
\end{align}
For the last inequality, the omitted harmonic sum is at most
$1+\log H_N\leq H_N/2$ when $H_N\geq8$.

\section*{One integrable terminal value}
For $m\geq1$, choose the explicit integers
\[
 N_m=\left\lceil\exp\left(8+a+
        \frac{160m^2(\log2)^2}{S\phi(1)^2}\right)\right\rceil,
 \qquad
 \xi=\sum_{m=1}^{\infty}2^{-m}X_{N_m}.
\]
Tonelli's theorem and \eqref{terminalN} give $\E\xi=1$, so this sum
is finite almost surely. The summands need not be independent.
Since $H_N\geq\log(N+1)$, \eqref{divergence} gives
$u^{N_m}_0\geq2^{2m}$.

Suppose a class-$(D)$ solution with terminal value $\xi$ existed.
For $0<c\leq1$, $g(cz)\geq c g(z)$, because the logarithmic denominator
decreases under scaling. Hence $2^{-m}u^{N_m}$ is still a bounded
nonnegative lower solution, with terminal value at most $\xi$.
Lemma~\ref{comparison} yields
\[
 Y_0\geq2^{-m}u^{N_m}_0\geq2^m\qquad(m\geq1).
\]
These countably many inequalities contradict the finiteness of $Y_0$.
This proves Theorem~\ref{main} and the negative answer to
\cite[Problem 5.1, equation (5.1)]{FHT}.

\paragraph{Scope and attribution.}
The generator, exponent, Brownian basis and class-$(D)$ condition are
unchanged. Only an admissible terminal value is chosen to refute the
universal existence claim. The proof and independent mathematical checks
were developed with Codex AI assistance. No Lean verification, external
peer review or publication priority is claimed.

\begin{thebibliography}{9}
\bibitem{FHT}
S. Fan, Y. Hu and S. Tang.
\emph{1D nonlinear backward stochastic differential equations:
a unified theory and applications}.
arXiv:2309.06233v2, 11 February 2026, \S1 and \S5, Problem 5.1.
\href{https://arxiv.org/html/2309.06233v2#S5}{Primary text}.
\end{thebibliography}
\end{document}
